
\section{A uniform nonlinear adaptive Gaussian theorem}
\label{sec:triangular}

This section defines the triangular experiment once and proves the probability
and Laplace estimates used later.  Its role is not to replace model-specific
identification.  The mechanical verification of the global contrast is given
in Proposition~\ref{prop:balanced-contrast}.

\subsection{Experiment class and two-stage filtration}

Let \(\Theta\subset\R^p\) be compact and convex, and let
\(\Theta_\circ\Subset\operatorname{int}\Theta\).  The hidden state space
\(\Hh\) is a separable Hilbert space.  Fix \(R<\infty\), noise variance
\(\sigma^2>0\), and an integer \(s>p/2\).

\begin{definition}[Triangular adaptive experiment]
\label{def:triangular}
For every terminal size \(n\), an admissible policy \(\kappa_n\) constructs
random variables in the chronological order
\[
 \cF_{i-1}\longrightarrow A_i\longrightarrow\xi_i\longrightarrow Y_i,
 \qquad i=1,\ldots,n.
\]
Here \(A_i\) is the action, possibly selected with a fresh policy randomizer,
\[
 \cG_i=\cF_{i-1}\vee\sigma(A_i),
 \qquad
 \cF_i=\cG_i\vee\sigma(Y_i),
\]
and \(\xi_i\sim N(0,\sigma^2)\) is independent of \(\cG_i\).  For every
\(\theta\in\Theta\), the zero-initial-state mean \(m_{n,i}(\theta)\) is a
\(\cG_i\)-measurable random \(C^{s+3}\) function of \(\theta\).  The
transient contribution of initial state \(z\) is
\(b_{n,i}(\theta,z)\), also \(\cG_i\)-measurable.  The external true law is
\begin{equation}
 Y_i=m_{n,i}(\theta_0)+b_{n,i}(\theta_0,z_0)+\xi_i,
 \qquad \theta_0\in\Theta_\circ,\quad\norm{z_0}\le R.
 \label{eq:true-triangular-law}
\end{equation}
The policy kernel is the same measurable function of the observed record
under every parameter; it cannot use \(\theta_0\) or \(z_0\) as an oracle.
It may depend on \(n\).
\end{definition}

For a working initial-state probability \(\nu_n\), supported on the radius
\(R\) ball but otherwise arbitrary, define
\begin{align}
 L_{n}^{\nu_n}(\theta)
 &=\int_{\Hh}\prod_{i=1}^n
 \varphi_\sigma\{Y_i-m_{n,i}(\theta)-b_{n,i}(\theta,z)\}\,\nu_n(dz),
 \label{eq:marginal-likelihood}\\
 \Pi_{n}^{\nu_n}(d\theta)
 &\propto \pi(\theta)L_n^{\nu_n}(\theta)\,d\theta .
 \label{eq:quasi-posterior}
\end{align}
Policy factors cancel in the chronological joint likelihood because the same
kernel is evaluated at the realized past.  If \(z_0\notin\supp\nu_n\),
\eqref{eq:quasi-posterior} is a quasi-posterior under the external truth
\eqref{eq:true-triangular-law}.  No density on \(\Hh\) is assumed.

The experiment index is
\[
 E=(n,\kappa_n,\theta_0,z_0,\nu_n)\in\frE_n.
\]
All random variables below are understood under
\(P_E=P_{\theta_0,z_0}^{\kappa_n}\).  The following notation makes every
uniform quantifier explicit.

\begin{definition}[Uniform stochastic orders]
\label{def:uniform-orders}
For random variables \(Z_E\), write \(Z_E=o_{\frE}(1)\) when, for every
\(\eta>0\),
\[
 \sup_{E\in\frE_n}P_E(\abs{Z_E}>\eta)\longrightarrow0.
\]
Write \(Z_E=O_{\frE}(1)\) when
\[
 \lim_{K\to\infty}\sup_{n\ge1}\sup_{E\in\frE_n}
 P_E(\abs{Z_E}>K)=0.
\]
The same definitions apply to random norms and suprema of random fields.
\end{definition}

\begin{assumption}[Uniform smoothness, transience, and contrast]
\label{ass:triangular}
There are deterministic constants \(B<\infty\), \(c_0>0\), a summable
sequence \((d_i)\), and \(\varepsilon_n\downarrow0\), independent of the
experiment index, such that:
\begin{enumerate}
\item for \(\abs\alpha\le s+3\),
\[
 \sup_{n,i,E,\theta}\abs{\partial_\theta^\alpha m_{n,i}(\theta)}\le B;
\]
\item for \(\abs\alpha\le s+3\),
\[
 \sup_{n,E,\theta,\norm z\le R}
 \abs{\partial_\theta^\alpha b_{n,i}(\theta,z)}\le d_i,
 \qquad \sum_{i\ge1}d_i<\infty;
\]
\item with probability at least \(1-\varepsilon_n\), simultaneously in
\(\theta\in\Theta\),
\begin{equation}
 \frac1n\sum_{i=1}^n
 \{m_{n,i}(\theta)-m_{n,i}(\theta_0)\}^2
 \ge c_0\abs{\theta-\theta_0}^2.
 \label{eq:abstract-contrast}
\end{equation}
\end{enumerate}
The prior density \(\pi\) is continuous on \(\Theta\), bounded above, and
bounded below on a fixed neighbourhood of \(\Theta_\circ\).
\end{assumption}

The summability hypothesis is stated with one deterministic envelope.  It
implies \(\abs{Y_i}\le B_0+\abs{\xi_i}\) for every \(i\) in the applications below;
equivalently this bound may be included in the experiment definition.

\subsection{The exact nuisance factor and misspecification}

Set
\[
 L_n^0(\theta)=\prod_{i=1}^n
 \varphi_\sigma\{Y_i-m_{n,i}(\theta)\}.
\]
Completing each Gaussian square gives the exact identity
\begin{align}
 L_n^{\nu_n}(\theta)&=L_n^0(\theta)e^{R_n^{\nu_n}(\theta)},
 \label{eq:nuisance-factorization}\\
 R_n^{\nu_n}(\theta)&=\log\int e^{V_n^\theta(z)}\,\nu_n(dz),\notag\\
 V_n^\theta(z)&=\sigma^{-2}\sum_{i=1}^n
 \left[(Y_i-m_{n,i}(\theta))b_{n,i}(\theta,z)
       -\tfrac12b_{n,i}(\theta,z)^2\right].
 \label{eq:nuisance-potential}
\end{align}

\begin{lemma}[Uniform nuisance derivative budget]
\label{lem:nuisance-budget}
For every fixed \(r\le s+3\) and \(q<\infty\), there is a random variable
\(W_{n,r,E}\) satisfying
\[
 \sup_{n,E}\E_E W_{n,r,E}^{\,q}<\infty
\]
and
\begin{equation}
 \sup_{\theta,\norm z\le R}
 \sum_{\abs\alpha\le r}
 \abs{\partial_\theta^\alpha V_n^\theta(z)}
 \le W_{n,r,E},
 \qquad
 \norm{R_n^{\nu_n}}_{C^r(\Theta)}
 \le C_r(1+W_{n,r,E})^{r+1}.
 \label{eq:nuisance-derivative-budget}
\end{equation}
The constants are independent of \(\nu_n\).
\end{lemma}

\begin{proof}
Every differentiated summand in \eqref{eq:nuisance-potential} is bounded by
\(C_rd_i(1+\abs{\xi_i})\).  Thus one may take
\[
 W_{n,r,E}=C_r\sum_{i\ge1}d_i(1+\abs{\xi_i}).
\]
Minkowski's inequality and the Gaussian moments prove the uniform moment
bound.  Derivatives of a log integral are finite sums of cumulants of the
bounded derivatives of \(V_n\) under the positive normalized tilt
\(e^{V_n}\nu_n/\int e^{V_n}d\nu_n\).  Their absolute values are bounded by a
universal polynomial in \(1+W_{n,r,E}\), proving
\eqref{eq:nuisance-derivative-budget} without any regularity assumption on
\(\nu_n\).
\end{proof}

To state the pseudo-true conclusion, let \(\theta_{n,E}^*\) be any minimizer
of the expected negative marginal log likelihood, normalized by \(n\).

\begin{proposition}[The transient pseudo-true parameter]
\label{prop:pseudotrue}
Suppose the failure probability in \eqref{eq:abstract-contrast} is
exponentially small.  Then
\[
 \sup_{E\in\frE_n}\abs{\theta_{n,E}^*-\theta_0}\le Cn^{-1}.
\]
In particular the asymptotic pseudo-true parameter is the physical parameter
\(\theta_0\), even when \(z_0\notin\supp\nu_n\).
\end{proposition}

\begin{proof}
The expected zero-state log-likelihood contrast is
\[
 \frac1{2\sigma^2n}\E_E\sum_i
 \{m_{n,i}(\theta)-m_{n,i}(\theta_0)\}^2
 -\frac1{\sigma^2n}\E_E\sum_i
 b_{n,i}(\theta_0,z_0)
 \{m_{n,i}(\theta)-m_{n,i}(\theta_0)\}.
\]
The first term is at least \(c\abs{\theta-\theta_0}^2\); the exponentially
small complement is harmless because all means are bounded.  The second is
bounded by \(C\abs{\theta-\theta_0}/n\).  By
Lemma~\ref{lem:nuisance-budget}, the expected difference
\(R_n^{\nu_n}(\theta)-R_n^{\nu_n}(\theta_0)\), divided by \(n\), has the same
linear bound.  A minimizer therefore obeys
\(c r^2\le Cr/n\), which proves the assertion.
\end{proof}

\subsection{Smooth martingale fields}

For a multi-index \(\beta\), put
\[
 S_{n,\beta}(\theta)=\sum_{i=1}^n
 \xi_i\partial_\theta^\beta m_{n,i}(\theta).
\]

\begin{lemma}[Policy-uniform Sobolev score bound]
\label{lem:sobolev-score}
For every \(\abs\beta\le3\),
\begin{equation}
 \E_E\norm{S_{n,\beta}}_{C(\Theta)}^2\le Cn
 \quad\hbox{uniformly in }E\in\frE_n.
 \label{eq:sobolev-score}
\end{equation}
Consequently \(n^{-1/2}\norm{S_{n,\beta}}_{C(\Theta)}=O_{\frE}(1)\).
\end{lemma}

\begin{proof}
Use the fixed bounded convex domain \(\operatorname{int}\Theta\), which has
the cone property, and its deterministic Sobolev embedding constant.
For each \(\abs\alpha\le s\), martingale orthogonality from
\(\E(\xi_i\mid\cG_i)=0\) gives
\[
 \E_E\int_\Theta
 \abs{\partial_\theta^\alpha S_{n,\beta}(\theta)}^2d\theta
 =\sigma^2\sum_{i=1}^n\E_E\int_\Theta
 \abs{\partial_\theta^{\alpha+\beta}m_{n,i}(\theta)}^2d\theta
 \le Cn.
\]
The deterministic Sobolev embedding \(H^s\hookrightarrow C\), valid because
\(s>p/2\), proves \eqref{eq:sobolev-score}.
\end{proof}

Write
\[
 g_i=Dm_{n,i}(\theta_0),\qquad
 I_n=\frac1{n\sigma^2}\sum_{i=1}^ng_ig_i^T,\qquad
 \Delta_n=\frac1{\sigma^2\sqrt n}\sum_{i=1}^ng_i\xi_i.
\]
The contrast event implies, by division along a line through \(\theta_0\),
\[
 I_n\ge(c_0/\sigma^2)\Id.
\]
The smoothness bound gives \(I_n\le(B^2p/\sigma^2)\Id\), and
Lemma~\ref{lem:sobolev-score} gives \(\Delta_n=O_{\frE}(1)\).

\subsection{Local quadratic expansion and global localization}

\begin{proposition}[Uniform local expansion and global envelope]
\label{prop:local-global}
For every fixed \(M\),
\begin{equation}
 \sup_{\abs h\le M}
 \left|
 \log\frac{L_n^{\nu_n}(\theta_0+h/\sqrt n)}
           {L_n^{\nu_n}(\theta_0)}
 -h^T\Delta_n+\tfrac12h^TI_nh
 \right|=o_{\frE}(1).
 \label{eq:uniform-lan}
\end{equation}
On the contrast event there is \(Z_{n,E}=O_{\frE}(1)\) such that,
simultaneously in \(\theta\in\Theta\),
\begin{equation}
 \log\frac{L_n^{\nu_n}(\theta)}{L_n^{\nu_n}(\theta_0)}
 \le-c_1n\abs{\theta-\theta_0}^2
       +Z_{n,E}\sqrt n\abs{\theta-\theta_0},
 \qquad c_1=\frac{c_0}{2\sigma^2}.
 \label{eq:global-envelope}
\end{equation}
Consequently, for every \(\eta>0\),
\begin{equation}
 \lim_{M\to\infty}\limsup_{n\to\infty}
 \sup_{E\in\frE_n}
 P_E\left\{
 \Pi_n^{\nu_n}(\sqrt n\abs{\theta-\theta_0}>M)>\eta
 \right\}=0.
 \label{eq:root-n-contraction}
\end{equation}
\end{proposition}

\begin{proof}
Put \(d_i(\theta)=m_{n,i}(\theta)-m_{n,i}(\theta_0)\).  The exact zero-state
log-likelihood ratio under the external truth is
\begin{equation}
 \sigma^{-2}\sum_i\xi_i d_i(\theta)
 +\sigma^{-2}\sum_i b_{n,i}(\theta_0,z_0)d_i(\theta)
 -\frac1{2\sigma^2}\sum_i d_i(\theta)^2.
 \label{eq:exact-zero-state-ratio}
\end{equation}
The fundamental theorem of calculus and
Lemma~\ref{lem:sobolev-score} bound its innovation term by
\(O_{\frE}(\sqrt n)\abs{\theta-\theta_0}\), uniformly in \(\theta\).
The true transient cross term is at most
\(C(\sum_{i\ge1}d_i)\abs{\theta-\theta_0}\), where the \(d_i\)'s inside
the summation are the deterministic transient envelope from
Assumption~\ref{ass:triangular}, not the mean increment \(d_i(\theta)\).
Lemma~\ref{lem:nuisance-budget} bounds the difference of the two \(R_n\)
terms by \(O_{\frE}(1)\abs{\theta-\theta_0}\). Combining with
\eqref{eq:abstract-contrast} proves \eqref{eq:global-envelope}.

For \(\delta=h/\sqrt n\), Taylor expansion gives
\[
 d_i(\theta_0+\delta)=g_i^T\delta
 +\tfrac12\delta^TD^2m_{n,i}(\theta_0)\delta+O(\abs\delta^3).
\]
The innovation multiplied by the Hessian is \(O_{\frE}(n^{-1/2})\) by the
score lemma, its third-order remainder is bounded by
\(Cn^{-3/2}\sum_i\abs{\xi_i}=O_{\frE}(n^{-1/2})\), and the quadratic-sum
remainder is \(O(n\abs\delta^3)=O(n^{-1/2})\) on a fixed \(h\)-ball. The
true transient and \(R_n\) differences are also \(O_{\frE}(n^{-1/2})\).
This proves \eqref{eq:uniform-lan}.

On events where \(Z_{n,E}\le K\), the rescaled right side of
\eqref{eq:global-envelope} is \(-c_1\abs h^2+K\abs h\). The local expansion
integrated over \(\abs h\le1\), together with the lower bound on \(\pi\),
gives a denominator at least \(n^{-p/2}e^{-O_{\frE}(1)}\). The Gaussian
integral of the global envelope outside \(\abs h>M\) tends to zero
uniformly as \(M\to\infty\). Tightness and the vanishing
contrast-failure probability prove \eqref{eq:root-n-contraction}.
\end{proof}

\subsection{A relative random-information Laplace lemma}

\begin{lemma}[Uniform relative Laplace principle]
\label{lem:random-laplace}
Let \(K_{n,E}\ge0\) on \(H_{n,E}=\sqrt n(\Theta-\theta_0)\).
Suppose, uniformly in the experiment class, that the eigenvalues of
\(J_{n,E}\) lie in \([\underline j,\overline j]\) with probability tending
to one, \(U_{n,E}=O_{\frE}(1)\), and for every fixed \(M\),
\[
 \sup_{|h|\le M}|\log K_{n,E}(h)-h^TU_{n,E}+\tfrac12h^TJ_{n,E}h|
 =o_{\frE}(1).
\]
Suppose also \(K_{n,E}(h)\le C_{n,E}e^{-a|h|^2+B_{n,E}|h|}\), with
\(C_{n,E},B_{n,E}=O_{\frE}(1)\), on these events. Let \(r_{n,E}\ge0\)
be uniformly bounded and converge uniformly to one on each fixed ball.
Extend \(K_{n,E}r_{n,E}\) by zero outside \(H_{n,E}\), and put
\(q_{n,E}(h)=e^{h^TU_{n,E}-h^TJ_{n,E}h/2}\). Then
\begin{equation}
 \frac{\int_{\R^p}|K_{n,E}r_{n,E}-q_{n,E}|dh}{\int_{\R^p}q_{n,E}dh}
 =o_{\frE}(1),\qquad
 \int q_{n,E}dh=(2\pi)^{p/2}(\det J_{n,E})^{-1/2}
 e^{U_{n,E}^TJ_{n,E}^{-1}U_{n,E}/2}.
 \label{eq:relative-L1}
\end{equation}
\end{lemma}
\begin{proof}
Restrict to the event \(|U|+B+C\le K\) and the stated spectral bounds.
The Gaussian integral, evaluated by completing the square, is bounded
above and below by positive constants depending only on \(K\) and the
fixed spectral interval. On \(|h|\le M\), exponentiation of the uniform
log error and the prior-ratio error gives vanishing \(L^1\) difference.
Outside this ball, both the Gaussian and the common quadratic envelope
have integrals tending uniformly to zero as \(M\to\infty\).
The scaled domains exhaust every fixed ball uniformly because the truth
lies in \(\Theta_\circ\). Divide by the lower Gaussian integral bound,
then let \(n\to\infty\), \(M\to\infty\), and \(K\to\infty\).
Relative \(L^1\) error controls the normalizers and hence total variation.
\end{proof}

\begin{theorem}[Uniform nonlinear adaptive quasi-Bernstein--von Mises]
\label{thm:abstract-bvm}
Under Assumption~\ref{ass:triangular}, uniformly over \(E\in\frE_n\),
\begin{equation}
 \TV\left(\mathcal L\{\sqrt n(\theta-\theta_0)\mid\cF_n\},
 N(I_n^{-1}\Delta_n,I_n^{-1})\right)=o_{\frE}(1).
 \label{eq:abstract-bvm}
\end{equation}
Define the comparison law arbitrarily when \(I_n\) is singular, an event
of uniformly vanishing probability. The evidence obeys
\begin{equation}
 \frac{\mathcal Z_n^{\nu_n}}{L_n^{\nu_n}(\theta_0)}
 =n^{-p/2}\pi(\theta_0)(2\pi)^{p/2}(\det I_n)^{-1/2}
 e^{\Delta_n^TI_n^{-1}\Delta_n/2}[1+o_{\frE}(1)].
 \label{eq:abstract-evidence}
\end{equation}
Neither assertion assumes convergence of \(I_n\).
\end{theorem}
\begin{proof}
Apply Lemma~\ref{lem:random-laplace} to the likelihood ratio at
\(\theta_0+h/\sqrt n\), with \(J=I_n\), \(U=\Delta_n\), and
\(r(h)=\pi(\theta_0+h/\sqrt n)/\pi(\theta_0)\).
Proposition~\ref{prop:local-global} supplies the local expansion and
Gaussian envelope; contrast and the score lemma supply the spectral and
tightness hypotheses. Prior continuity and positivity give the ratio
condition. Normalize the relative \(L^1\) bound to obtain total variation,
and retain its integral with the Jacobian \(n^{-p/2}\) to obtain evidence.
\end{proof}
